\documentclass[10pt]{article}
\usepackage[landscape,margin=0.8in]{geometry}
\usepackage{amsmath,amssymb,booktabs,longtable,array}
\newcommand{\Hist}{\mathcal{H}}
\newcommand{\rep}{\sigma}
\newcommand{\NN}{\mathbb{N}}

\title{Proof Spine for \emph{Distinction and Continuation} and \emph{The Shape of Possibility}}
\author{Flyxion\\Independent Researcher}
\date{Working ledger, October 2026}

\begin{document}
\maketitle

\noindent Each load-bearing result of the two volumes appears below with its status. \textbf{E} marks an established result, to be attributed to its source; \textbf{N} marks a new result to be proved in full under the listed assumptions; \textbf{O} marks a proof obligation. A result with status O may not be cited as a theorem by any later chapter. Volume~1 is \emph{Distinction and Continuation} and Volume~2 is \emph{The Shape of Possibility}.

\small
\begin{longtable}{@{}l p{0.30\textwidth} c l p{0.22\textwidth} p{0.14\textwidth}@{}}
\toprule
\textbf{ID} & \textbf{Statement} & \textbf{St.} & \textbf{Where} & \textbf{Assumptions / source} & \textbf{Depends on}\\
\midrule
\endhead
S1 & A decision factors through a representation iff equal representations force equal decisions. & E & V1 Ch.~2; both Ch.~1 & Standard quotient factorization. & none\\[3pt]
S2 & If $\rep(h_1)=\rep(h_2)$ and $P(h_1)\ne P(h_2)$, no procedure receiving $\rep$ alone is correct on both histories. & E & both Ch.~1 & Corollary of S1. & S1\\[3pt]
S3 & Under a sound gate that consults $\rep$ alone, a continuation admissible at $h_1$ is refused whenever some $h_2$ with the same representation is inadmissible. & N & both Ch.~1 & Soundness over all of $\Hist$; gate depends on $\rep(h)$ only. Elementary. & S1\\[3pt]
S4 & Representational obstruction can be inherited from an obligation obstruction at a different history. & N & both Ch.~1 & Example, not a theorem. & S3\\[3pt]
S5 & Coarsest representation deciding a fixed policy is the kernel of the policy. & E & V1 Ch.~5 & Standard. & S1\\[3pt]
S6 & For a fully separating policy family, every sufficient representation is injective. & N & V1 Ch.~6 & Separating means distinct histories are distinguished by some policy in the family. Elementary. & S5\\[3pt]
S7 & Minimal automaton and continuation equivalence (Myhill--Nerode). & E & V1 Ch.~7 & Nerode (1958); Myhill (1957). & none\\[3pt]
S8 & Gap between one-shot and update-maintainable sufficiency, characterized by the Myhill--Nerode congruence over the update family. & N & V1 Ch.~7 & $s_{\mathrm{one}}$ finite for the equality criterion; counter-families $P_n$ and $q_{\mathrm{neg}}$ given. Proved in full. & S7\\[3pt]
S9 & Tight state complexity $mn$ for the intersection of two independent modular counters over disjoint alphabets. & E & V1 Ch.~8 & Yu, Zhuang, Salomaa (1994); generalized in S36. & S7\\[3pt]
S10 & Admissibility decomposes into physical executability, authorization, and obligation satisfaction; the normative obstruction has four forms, one of them joint. & N & V1 Ch.~13 & Prefix-closed sets; definitional plus elementary. & none\\[3pt]
S11 & Monitor for bounded-response obligations over disjoint alphabets: the minimal monitor has exactly $\prod_i(k_i+1)+1$ states. & N & V1 Ch.~14 & Fragment fixed (Definition of the disjoint fragment), counting in each component's own events. Shared-clock case remains open. Proved in full. & S7, S36\\[3pt]
S12 & Coarsest identity criterion preserving promised observations is their joint kernel; preserving a promise with updates (a submonoid) it is the continuation equivalence of the family. & N & V1 Ch.~15 & Elementary; $F$ a submonoid. Proved in full. & S5, S35\\[3pt]
S13 & Admissible shedding characterized from dependency structure (disjoint from a support); no additive capacity measure suffices when alternatives exist. & N & V1 Ch.~17 & Finite Horn dependency system and required set given; proved in full (see S42). & S10\\[3pt]
S14 & Coarsening a representation contracts the family of implementable policies; later computation cannot recover discarded distinctions. & N & V1 Ch.~21 & $\rep'=f\circ\rep$. Elementary, via S1. & S1\\[3pt]
S15 & Belief state is a sufficient statistic in a POMDP. & E & V2 Ch.~5 & \r{A}str\"om; Smallwood and Sondik. & none\\[3pt]
S16 & Belief state is minimal sufficient for a stated policy family. & O & V2 Ch.~5 & Policy family and observation model to be fixed. & S15, S5\\[3pt]
S17 & Augmented state sufficient for feasibility iff the added component is sufficient in the V1 sense. & N & V2 Ch.~7 & Feasibility is a function of the history. Elementary. & S5\\[3pt]
S18 & Discrete-time viability kernel is the greatest fixed point of the one-step predecessor operator. & E & V2 Ch.~3 & Aubin; standard. & none\\[3pt]
S19 & Reachability decidable for Petri nets; coverability EXPSPACE-hard. & E & V2 Ch.~6 & Mayr; Kosaraju; Lipton. & none\\[3pt]
S20 & Viability kernels for systems whose constraint set is updated by an obligation monitor. & N & V2 Ch.~14 & Bounded-response monitor of S11 in product form; kernel monotone in slack; corridor example in full. & S11, S18\\[3pt]
S21 & Invariant preservation for explicit/automatic switching under dwell time and monitoring latency. & N & V2 Ch.~17 & Each controller keeps the safe set invariant alone; latency bound; failure model. Liberzon; Hespanha and Morse. & S18\\[3pt]
S22 & Interface preserves admissible reachability iff a restricted simulation relation exists. & N & V2 Ch.~18 & Admissible region fixed; abstraction total on it. & S10\\[3pt]
S23 & Invariance of a constraint set under rule changes restricted to an explicit relation, including constraints on revising constraints. & N & V2 Ch.~19 & Rule-change relation and its own constraint stated. & S18\\[3pt]
S24 & Proxy-optimal policy is $2\varepsilon/(1-\gamma)$-optimal under uniform reward error $\varepsilon$. & E & V2 Ch.~20 & Discounted MDP; Puterman. & none\\[3pt]
S25 & Constrained-MDP version of S24 with proxied objective and constraints. & N & V2 Ch.~20 & Uniform bound on both objective and constraint error; feasibility may fail. & S24\\[3pt]
S26 & Cost bounds for maintaining a representation sufficient for an anticipated policy family. & N & V2 Ch.~22 & Cost model to be fixed. & S14, S6\\[3pt]
S27 & Continuation quotient of the account: classes are $\bot$ (overdrawn) and the balances $n\in\NN$ otherwise. & N & V1 Ch.~2 & Positive integer amounts; overdraft permitted; all finite event sequences. Proved in full. & S1\\[3pt]
S28 & Under a size bound $2\le k<n$ there is no most informative representation of an $n$-element set, and at least two incomparable maximal ones. & N & V1 Ch.~3 & $X$ finite; size is number of values. Proved in full. & S1\\[3pt]
S29 & Myhill--Nerode construction for properties of histories: continuation equivalence is the coarsest right congruence inside $\ker P$, and the Nerode system is minimal. & E & V1 Ch.~4 & Myhill; Nerode. & S1\\[3pt]
S30 & Every behavioural property is the intersection of a safety and a liveness property. & E & V1 Ch.~4 & Alpern and Schneider (1985). & none\\[3pt]
S31 & Nonempty safety properties correspond bijectively to pruned history sets via the viable histories. & E & V1 Ch.~4 & Closed sets and pruned trees; Kechris. & S30\\[3pt]
S32 & On viable histories of the no-overdraft property, the set of admissible extensions determines the balance. & N & V1 Ch.~4 & Positive integer amounts. Proved in full. & S27, S31\\[3pt]
S33 & Sufficiency for a fixed policy: $\ker r\subseteq\ker\pi$; no non-injective representation suffices for all two-action policies; minimal representations are sensitive to a one-case change; sufficiency depends on the domain. & E, N & V1 Ch.~5 & Theorem E; the other three N, elementary, proved in full. & S1\\[3pt]
S34 & Joint sufficiency for policy families; separating families force injectivity; threshold chains cost $K+1$ values. & E, N & V1 Ch.~6 & Bounds and VC bound E (Sauer--Shelah, Assouad); separating and chain results N. Sharp bounds remain a research problem. & S5, S6\\[3pt]
S35 & Maintainable and $F$-maintainable sufficiency; Nerode theorem for update families; minimal machine by Moore partition refinement. & E, N & V1 Ch.~7 & Moore (1956), Hopcroft (1971) cited; $F$-congruence theorem N. & S7, S8\\[3pt]
S36 & Independent properties over disjoint alphabets: minimal memory is the product of the classes; sufficiency composes but minimality does not. & N & V1 Ch.~8 & Alphabet partition with each property depending on its own part. Proved in full. & S7, S9\\[3pt]
S37 & Exact sufficiency is the only preservation holding under every distribution; Bayes error is monotone under coarsening and errors add across policies (tight); zero conditional entropy characterizes sufficiency. & E, N & V1 Ch.~9 & Fano and data-processing inequalities E. Others N, finite $X$, proved in full. & S5\\[3pt]
S38 & Verified region of a claim is its interior in the partition induced by the evidence; a proof reaches a case only through evidence settling its assumptions. & E, N & V1 Ch.~10 & Certifying algorithms E (McConnell et al.; Necula). Interior and gap results N, proved in full. & S1\\[3pt]
S39 & Hash chain determines the sequence (idealized injective hash); vector clocks; provenance semirings. Minimal last-author record has $(|U|+1)^K$ values; content identity does not determine lineage. & E, N & V1 Ch.~11 & First three E (Merkle; Fidge, Mattern; Green et al.). Counts and separation N, proved in full. & S6, S36\\[3pt]
S40 & Parity needs $n$ probes and a marked cell $n-1$; derived cells leave storage adequacy unchanged while collapsing retrieval depth. & N & V1 Ch.~12 & Deterministic cell-probe readers; worst case over the domain. Proved in full. & S1\\[3pt]
S41 & Exactly-once effects under identity, atomic record, and fair channels; non-atomic variants give at-most-once or at-least-once; deduplication memory is exactly $2^{|I|}$ states under unbounded duplication. & E, N & V1 Ch.~16 & Two generals impossibility E (Akkoyunlu et al.; Gray; Halpern and Moses). Rest N, proved in full. & S7\\[3pt]
S42 & Admissible shed sets are those disjoint from a self-supporting set; scalar capacity tests are correct for conjunctive dependencies and fail with alternatives (counterexample). Smallest support is set-cover hard. & N & V1 Ch.~17 & Finite Horn dependency systems. Proved in full; hardness by standard reduction. & S10, S11\\[3pt]
S43 & Restoration and replacement defined relative to an identity criterion; no compensation restores a step crossing a bad prefix; aborted sagas restore iff criterion is a congruence for compensations (counterexample otherwise). & N & V1 Ch.~18 & Prefix-closed obligation sets; per-step restoration. Proved in full. & S12, S30\\[3pt]
S44 & Global property decidable from causal past iff insensitive to concurrent events; quota theorem: independent capacities of replicas sum to at most the common balance. & N & V1 Ch.~19 & Finite executions, integer withdrawals. CRDT convergence and invariant confluence E (Shapiro et al.; Bailis et al.). Proved in full. & S39\\[3pt]
S45 & Refinement preserves compatibility of deterministic interface automata; non-circular assume--guarantee rule under alphabet conditions; attenuating delegation creates no authority. & N & V1 Ch.~20 & Deterministic automata; prefix-closed components. Nondeterministic version E (de Alfaro, Henzinger). Proved in full. & S10\\[3pt]
S46 & Revision horizon determines the representation and shrinks by $|A|^{-(k-k')}$ under coarsening; affected cases are the non-constant classes; worst-case binary revision error under the uniform distribution is $\frac1n\sum_y\lfloor n_y/2\rfloor\ge\frac{n-k}{2n}$, against $\approx1/2\varepsilon$ values for thresholds. & N & V1 Ch.~21 & Finite $X$, uniform $\mu$. Proved in full; general approximate revision sufficiency open. & S1, S14, S37\\[3pt]
S47 & Representation selection is weighted set cover over disagreement pairs: greedy $H(|D|)$ approximation, NP-hard (injective policy is minimum test set), checking is $O(|X||F|)$; updatable storage for $P_n$ costs exactly $n+1$ values. & E, N & V1 Ch.~22 & Minimum test set and greedy cover E (Garey--Johnson; Chv\'atal). Rest N, proved in full. & S8, S35\\[3pt]
S48 & Worked systems: 13-state monitor and gate optimality; document record adequacy; checkout service supports; ledger semantics. & N & V1 Ch.~23 & Each instantiates S11, S13, S39. Proved in full. & S11, S39, S42\\
S49 & SCC partition is the minimal sufficient representation for the reachability family; irreversible edges are exactly the inter-component edges; at most $c-1$ component changes. & N & V2 Ch.~2 & Finite transition graph. Tarjan cited. Proved in full. & S1\\[3pt]
S50 & Greatest fixed point of the predecessor operator is the viability kernel and the set of states with infinite $K$-trajectories; viable histories are those ending in the kernel; kernel convex for linear systems with convex $U,K$. & N & V2 Ch.~3 & Finite branching or compactness as stated. Aubin cited. & S18\\[3pt]
S51 & Almost-sure safety is the support-based viability kernel; optimizing under it reduces to the restricted MDP; expected-cost constraints can force randomization. & N & V2 Ch.~4 & Finite MDP. Bellman contraction and occupation measures cited. & S18\\[3pt]
S52 & Predictive equivalence of beliefs is the orthogonal complement of the span of prediction vectors; belief minimal iff rank condition; almost-sure safety decided on supports. & N & V2 Ch.~5 & Finite POMDP. Proved in full; refines S15. & S15\\[3pt]
S53 & State equation necessary not sufficient for Petri net reachability; place invariants bound markings; the viability set of a net is upward closed with finite basis. & N & V2 Ch.~6 & Nonnegative integer weights. Dickson cited. & S19\\[3pt]
S54 & Augmentation theorem: least augmentation has size equal to the fibre count; hard budgets are exact by the remaining-budget product; unbounded counters admit no finite augmentation. & N & V2 Ch.~7 & Maintainable state description. Completes S17. & S17, S29\\[3pt]
S55 & Return cost is a quasi-metric whose finite part is reversibility; budget above $(|X|-1)c_{\max}$ buys nothing; path dependence is insufficiency of the state; policies refine components. & N & V2 Ch.~8 & Finite graph, nonnegative costs. Norris cited for absorption. & S49, S54\\[3pt]
S56 & In monotone capability systems reachable configurations have a unique maximum, the closure, computed in linear time; shortest acquisition plans have length at most $|\mathrm{Cap}\setminus C|$; nominal, eventual, current menus are nested. & N & V2 Ch.~9 & Monotone (no losses). Bylander cited for the non-monotone case. & S49\\[3pt]
S57 & Acquisition with costly inspection: backward induction on a finite sufficient state, oracle bound with explicit gap, submodularity of domain value, budget-admissible optimum. & N & V2 Ch.~10 & Independent finite-support values for the DP; submodularity needs only nonnegativity. Weitzman cited. & S54\\[3pt]
S58 & Cardinality, volume, and graph distance of reach misrank budgeted freedom; dominance is a preorder with incomparable pairs and weighted scores reverse them; simulation preserves coverage; growth rate of feasible histories is the largest spectral radius among reachable components. & N & V2 Ch.~11 & Finite graph. Milner; Perron--Frobenius cited. & S49\\[3pt]
S59 & Cost algebras change the optimal path; componentwise optima not jointly achievable; scalar weights decide convex mixtures; Pareto front can have $2^n$ points; integer vector budgets decided by product; shared acquisition subadditive; time and energy of task graphs. & N & V2 Ch.~12 & Finite graphs, nonnegative costs. Lehmann, Garey--Johnson, Graham cited. & S55, S56\\[3pt]
S60 & Robust kernel is the greatest fixed point for arbitrary spaces; monotone in disturbance set; sharp threshold $\delta=\tfrac12$ in the scalar example; Gaussian chance-constraint equivalent; Boole allocation; no forever-safety under unbounded noise; constant parameters can be strictly easier than varying ones. & N & V2 Ch.~13 & History-dependent policies; bounded or full-support noise as stated. & S18, S51\\[3pt]
S61 & Compliant viability is the kernel of the product with the obligation monitor; monotone in slack; a trigger is admissible iff its successor is in the kernel; explicit corridor kernel; trigger admissibility depends on pending obligations; no sound and complete monitor-only gate. & N & V2 Ch.~14 & Finite deterministic system, disjoint trigger and response sets. Completes S20. & S11, S18\\[3pt]
S62 & Viable optimization is optimization on the kernel; price of viability zero iff an unconstrained optimum stays in $K$ and unbounded in general; lookahead $|X|$ needed and sufficient; $m$-option kernels nested, branching $m^k$, prices nondecreasing in $m$. & N & V2 Ch.~15 & Finite deterministic discounted system. & S18, S51\\[3pt]
S63 & Recoverable region is the least fixed point of the backward step and predecessor closed; progress functions bound recovery time by $\varphi$ and greedy descent with any fixed lookahead can leave the region; attractor strictly smaller than possible recovery; fallback regions of shedding are down-closed, linear for a unique support, NP-complete in general, exponentially many maximal elements. & N & V2 Ch.~16 & Finite systems; dependency systems of S42. & S42, S49\\[3pt]
S64 & Hybrid switching: a common invariant set needs no switching conditions; distinct invariant sets do not compose; with a common core, handover time and latency the safe set is preserved under a dwell-time clause; sharp example; tolerable latency equals exit distance. & N & V2 Ch.~17 & Finite or compact systems; fault model for the automatic controller. & S21\\[3pt]
S65 & Interfaces as macro denotations: admissible reach is characterized by weak simulation in both directions (completeness and soundness); one-step sufficiency is not necessary; forced reachability of the presented system lies in possible recovery of the concrete one. & N & V2 Ch.~18 & Admissible step relations fixed. & S22\\[3pt]
S66 & Augmented rule systems: the robust kernel has sections that are viability domains relative to the intersection of amendment sections; the amendment depth is a shortest path in the revision graph; unrestricted revision gives depth one; every depth is realized. & N & V2 Ch.~19 & Finite rule and state sets; amender acts after the controller. & S23, S60\\[3pt]
S67 & Proxy transfer is tight at $2\varepsilon/(1-\gamma)$; on-policy refinement charges error on the two occupancy measures; constrained version with tightened budget has loss $2\rho$ plus a concavity-bounded price; hard safety on a viability domain is independent of the proxy. & N & V2 Ch.~20 & Finite MDPs; discounted criterion. & S24, S25\\[3pt]
S68 & Untried versus infeasible: trichotomy of attempted, certified infeasible and untested; data on attempted alternatives are silent on independent untested ones; threshold horizon $n^*=1+c/(p\Delta)$ and minimax regret; knowledge-indexed viability kernel is monotone in knowledge and safe attempts are those with all outcomes inside it. & N & V2 Ch.~21 & Static feasibility; revealed outcomes deterministic. & S18, S12\\[3pt]
S69 & Least maintainable representation sufficient for a family is the intersection of Nerode congruences, with between the largest and the product of class counts; marginal cost of one more property; rent-or-buy with ratio 2 optimal among deterministic rules; irreversible loss forces indefinite maintenance; shared rent. & N & V2 Ch.~22 & Cost model of Def.~22.2; exponential demand for the stochastic parts. & S26, S14\\[3pt]
S70 & Flexibility-constrained design: commit iff $g>hR$ independently of horizon; polynomial for single-route interventions, NP-complete with two routes (vertex cover); flexibility price is the gain of the union of routes and decreases under refinement. & N & V2 Ch.~23 & Additive gains; independent exponential demand. & S66, S69\\[3pt]
S71 & Four case studies: vehicle kernel $b\ge x$ and robust $b\ge(1+\delta)x$; archive amendment depth and 8-value representation; adaptive controller kernels $K$, $\{-3..1\}$, $\varnothing$; network fallback region, redundancy is not capacity. & N & V2 Ch.~24 & Finite instances computed exactly. & S66, S69, S67\\[3pt]
S72 & Appendices B to G (self-contained recaps): walk counts grow at the spectral radius; linear-time kernel algorithm with counters; scalar affine invariant intervals; Lyapunov recovery time; concavity of LP value in the budget; value-iteration stopping rule; memoryless laws; Monte Carlo sample size. & N & V2 App.~B--G & Standard background with short proofs; cited results marked established. & S18, S58\\[3pt]
S73 & Book 1 appendices A to F: induction on histories, finite Knaster--Tarski, fooling sets, adversary and exchange arguments; representations form a partition poset with products and coproducts, quotients as coequalizers, Galois connection of questions and representations; Moore minimization and $(k+2)$-state bounded-response monitor; data processing and sufficiency, non-recovery bound; self-contained recap of reachability and the viability kernel with a linear-time algorithm. & N & V1 App.~A--F & Finite sets throughout; classical results marked established. & S18, S29\\[3pt]
S74 & Expansion of Chapters 1 to 10: representation without the gate bit permits no sample; reach profile stabilizes in at most $n-1$ steps; DAG of components and path counts; kernel algebra and stopping curve; geometric ruin and one-random-state bound; down-closed winning supports; backward coverability and Karp--Miller; decidability of sufficiency and normal-form rate-limit states; Eswaran--Tarjan augmentation and expected return cost; depth versus plan length and non-submodular closure; spread raises the reservation value (index policy $8.4$ against $5.5$ for the mean rule). & N & V2 Ch.~1--10 & Finite systems; instances computed by script. & S10, S18, S31\\[3pt]
S75 & Dominance is the componentwise order on cost vectors; simulation is monotone for label growth; rate-limit growth rates $\varphi$ and $1.4656$ from Fibonacci recurrences. & N & V2 Ch.~11 & Finite graphs; deterministic for $L_k=N_k$. & S40\\[3pt]
S76 & Lagrangian bound equals the mixture optimum (gap $6$ against $4$ on a three-path mission); list-scheduling anomaly ($12$ with three workers, $15$ with four). & N & V2 Ch.~12 & Finite path set; Graham list rule. & S41\\[3pt]
S77 & Distribution-free chance margin $\sqrt{(1-\varepsilon)/\varepsilon}$ against Gaussian margin; horizon allocation numbers; asymptotic geometric exit time via Krein--Rutman, with survival table. & N/E & V2 Ch.~13 & Strictly positive Gaussian kernel on a compact set. & S42, S43\\[3pt]
S77a & Exact admissibility of a second trigger in the corridor (two orderings); corridor kernel count $(L+1)+\sum_{c\le k}\min(c,L+1)$. & N & V2 Ch.~14 & Responses at the two ends; verified on $2646$ cases. & S44\\[3pt]
S78 & The unit-control speed-limited line has $W_2=\varnothing$; doubling authority gives kernels of sizes $77,59,0$; fault tolerance is a transversal number minus one (navigation stack, $|\mathrm{FB}(3)|=20$, $|\mathrm{FB}(4)|=31$). & N & V2 Ch.~15--16 & Finite systems; counts by exhaustive computation. & S45, S46\\[3pt]
S79 & Exact safe region $\{\ell\le d,\ \tau\ge\ell\}$ for relegation in the sharp family; minimum interface and size--effort trade-off; depth of a quorum rule equals the sum of the $q$ cheapest capture costs. & N & V2 Ch.~17--19 & Families as stated; game solved exhaustively for four parameter pairs. & S47, S48, S49\\[3pt]
S80 & Selection on a noisy Gaussian proxy delivers $m_n/\sqrt{1+\sigma^2}$ and the proxy overstates by exactly $1+\sigma^2$. & N & V2 Ch.~20 & Independent Gaussian values and noise. & S50\\[3pt]
S81 & Threshold for attempting once in a Bayesian two-armed instance, and the exact optimal values by dynamic programming; counters on one event need $\operatorname{lcm}$ classes; randomized rent-or-buy ratio $1.5353$ at $B=10$. & N/E & V2 Ch.~21--22 & Uniform prior; one event; discrete rent-or-buy. & S51, S52\\[3pt]
S82 & Flexibility price equals minimum weight vertex cover (weighted instance: loss $14$ optimal, matching heuristic $29$, local ratio $16$); chance kernel of the vehicle is $b\ge x+\mathrm{Bin}$-quantile. & N/E & V2 Ch.~23--24 & Single-module routes; independent two-valued costs. & S53, S54\\[3pt]
S83 & Book 1 Chapters 1 to 4: the document record takes $(2^{k-2}(k+1))^m$ values; class counts $t+1$, $t+1$, $2t+1$ for ever, now and the pair in the monitor quotient; maximal elements of the adequacy order $R_k$ are the $k$-block partitions, counted by Stirling numbers ($S(6,3)=90$ of $203$); viability of a prefix-closed set equals viability of its closure, and the extension clause of continuation equivalence cannot be dropped (three Nerode states). & N & V1 Ch.~1--4 & Finite histories; counts confirmed by script. & S7\\[3pt]
S84 & Book 1 Chapters 5 to 8: a $k$-valued representation serves at most $k-1$ thresholds and is sufficient for $|A|^k$ policies; the marginal-policy count $\sum_c|\pi(c)|$ (four to six classes) with an XOR counterexample to monotone marginal cost; the single-use rate limit over $\mathrm{Hist}_w$ has $w$ classes; a product holds under shared events when each state is reachable by its own letters (counts $12$, $4$, $4$). & N & V1 Ch.~5--8 & Finite sets; counts by script. & S5, S9\\[3pt]
S85 & Book 1 Chapters 9 to 12: binned error table for the geometric law ($q=0.9$, $B=8$) and worst-case Bayes error $1-1/m$; interior of ``never overdrawn'' under the balance map ($4,7,13,23,43$ states); recency list of $\sum_j n!/(n-j)!$ values ($5,16,65$ against $5,13,33$); depth-one packing needs $\lceil m/\lfloor\log_2 s\rfloor\rceil$ derived cells. & N & V1 Ch.~9--12 & Finite supports; counts by script. & S5, S8\\[3pt]
S86 & Book 1 Chapters 13 to 16: joint obstructions number $3^r-2^{r+1}+1$ ($50$ of $81$ at $r=4$); a minimal conflict can have any size $k$ (Sperner bound); two obligations on a shared clock with equal deadline $k$ have $k^2+k+1$ live reachable states (proved), and the minimal realized monitor has $k^2+k+2$ states (proved for every pair of deadlines: no two live states merge); one-level undo gives $k^2$ identity classes, redo $6$, depth-two stack $8$; apply-then-record gives a geometric number of applications with mean $1/(1-c)$. & N & V1 Ch.~13--16 & Finite instances; the three-obligation count remains computational. & S10, S11, S12\\[3pt]
S87 & Book 1 Chapters 17 to 20: admissible shed sets number $4^m-3^m$ with best scalar tests as computed (the integer-weight bound for $m=3$ is an upper bound only); under the identity criterion $\approx^F$ with the compensations in $F$ the congruence condition is automatic; a deposit restores a withdrawal iff $q_{\mathrm{neg}}(h)=1$ or $x\le\beta(h)$; fixed quotas serve at worst $1/n$ of the demand; the fewest delegations to withhold equals the maximum number of edge-disjoint chains (Menger). & N/E & V1 Ch.~17--20 & Finite systems; Menger's theorem cited. & S13, S12\\[3pt]
S88 & Book 1 Chapters 21 to 23: $k+1$ of the $n+1$ threshold revisions are served exactly by $k$ intervals, and $\lceil n/(2\lfloor\varepsilon n\rfloor+1)\rceil$ intervals give worst error at most $\varepsilon$; parity over a threshold vocabulary requires all $n-1$ observations; the checkout service has $6$ admissible shed sets of $16$. & N & V1 Ch.~21--23 & Uniform law; interval partitions. & S14, S13\\[3pt]
S89 & Beyond prediction, the belief is minimal sufficient for the observation probabilities and the values of all finite-horizon policies exactly when $(B-B)\cap Q^{*\perp}=\{0\}$; $V_n$ depends on the belief only through its projection onto $Q^{*}$; the free-parameter count for prediction is $\dim Q-1$. & N & V2 Ch.~5 & Finite POMDP; discount in $(0,1]$. & S50\\[3pt]
S90 & Reachability with deletions equals monotone closure when no relevant capability is lost; the relaxation is a certificate of infeasibility in general; one relevant loss defeats the equivalence (counterexample on four capabilities). & N & V2 Ch.~9 & Gain takes precedence over loss. & S10\\[3pt]
S91 & Policies sufficient for two representations number $|A|^c$ with $c$ the classes of the join of the kernels; partition refinement stabilizes in at most $N-1$ rounds; the viable states of a finite structure are the stable set of a decreasing sequence with at most $|Q|$ steps; three counterexamples isolating maintainability, finiteness of the alphabet and one-shot versus maintained size. & N/E & V2 App.~A & Finite structures. & S1\\[3pt]
S92 & Two obligations on a shared clock: no two live reachable states merge, so the minimal realized monitor has one more state than there are live reachable states ($k^2+k+2$ for equal deadlines $k$). & N & V1 Ch.~14 & Two obligations; shared alphabet. & S86\\[3pt]
\bottomrule
\end{longtable}

\end{document}
